\documentclass[11pt,a4paper,oneside]{report}

\usepackage[utf8]{inputenc}
\usepackage[T1]{fontenc}
\usepackage[french]{babel}
\usepackage{lmodern}
\usepackage{geometry}
\usepackage{microtype}
\usepackage{graphicx}
\usepackage{array}
\usepackage{booktabs}
\usepackage{tabularx}
\usepackage{enumitem}
\usepackage{pdflscape}
\usepackage{xcolor}
\usepackage{tikz}
\usepackage{fancyhdr}
\usepackage{hyperref}
\usetikzlibrary{arrows.meta,positioning,shapes.geometric,calc}

\geometry{margin=2.4cm}
\setlength{\headheight}{14pt}
\setlist[itemize]{leftmargin=1.5em,itemsep=0.3em,topsep=0.4em}
\definecolor{fleetgreen}{HTML}{215B3C}
\definecolor{fleetblue}{HTML}{315B76}
\definecolor{fleetlight}{HTML}{EFF5F1}
\definecolor{fleetgray}{HTML}{58645D}
\hypersetup{
  colorlinks=true,
  linkcolor=fleetgreen,
  urlcolor=fleetblue,
  pdftitle={Rapport de stage -- Fleet Tracker -- Version 1.0.0},
  pdfsubject={Présentation fonctionnelle et architecturale de Fleet Tracker}
}

\pagestyle{fancy}
\fancyhf{}
\lhead{\small Fleet Tracker}
\rhead{\small Rapport de stage -- Version 1.0.0}
\cfoot{\thepage}
\renewcommand{\headrulewidth}{0.4pt}

\newcommand{\important}[1]{\textcolor{fleetgreen}{\textbf{#1}}}
\newcommand{\entity}[1]{\textbf{#1}}

\begin{document}

\begin{titlepage}
  \centering
  \vspace*{2.2cm}
  {\Large\scshape Rapport de stage\par}
  \vspace{0.7cm}
  {\Huge\bfseries\color{fleetgreen} Fleet Tracker\par}
  \vspace{0.5cm}
  {\Large Système de suivi de flotte de taxis\par}
  \vspace{1.3cm}
  {\large Présentation du besoin, de la solution et de son architecture\par}
  \vfill
  \begin{tikzpicture}
    \draw[fleetgreen,line width=1.5pt] (0,0) -- (10,0);
    \fill[fleetgreen] (5,0) circle (0.12);
  \end{tikzpicture}
  \vfill
  {\large Version 1.0.0\par}
  \vspace{0.25cm}
  {\large 9 octobre 2026\par}
  \vspace*{1.4cm}
\end{titlepage}

\pagenumbering{roman}
\tableofcontents
\clearpage
\pagenumbering{arabic}

\chapter*{Résumé}
\addcontentsline{toc}{chapter}{Résumé}

Fleet Tracker est une solution de suivi de taxis conçue pour donner à un
responsable de flotte une vue actualisée de ses chauffeurs et de leurs
positions. Le chauffeur partage sa localisation depuis une application mobile ;
un tableau de bord web présente les conducteurs sur une carte et signale leur
état de disponibilité. Une console distincte permet à un responsable de la
plateforme de consulter et de créer des centres.

Le système repose sur trois interfaces clientes, une API centrale et une base
de données relationnelle. Le projet constitue un socle de suivi opérationnel
pour une première mise en œuvre en Tunisie, avec un pilote envisagé dans le
Grand Tunis. Il ne comprend pas encore la répartition des courses, la gestion
des commandes ou les fonctions d'un progiciel de gestion intégré.

\chapter{Introduction et contexte}

\section{Présentation du projet}

Le secteur du taxi s'appuie sur des chauffeurs mobiles et des responsables qui
doivent pouvoir vérifier la disponibilité des véhicules et leur position. Le
projet Fleet Tracker répond à ce besoin au moyen d'une application chauffeur,
d'un tableau de bord d'administration et d'un service central de traitement.
Les différentes interfaces communiquent avec ce service plutôt que de partager
directement leurs données entre elles.

La solution est pensée pour plusieurs centres. Les chauffeurs et les
administrateurs sont rattachés à un centre, ce qui permet à chaque responsable
de consulter les données de son périmètre. Une console de plateforme distincte
est destinée à la gestion transversale des centres et à la consultation du
journal des événements de cette console.

\section{Problématique}

Sans outil commun de géolocalisation, le suivi d'une flotte repose facilement
sur des échanges ponctuels et des informations qui ne sont pas visibles au
même endroit. Le responsable peut alors manquer de visibilité sur les
chauffeurs disponibles, leur dernière position connue ou l'actualité de
l'information reçue.

\medskip
\noindent\textbf{Problématique retenue :} comment fournir au responsable d'un
centre une vision simple et suffisamment récente de l'activité et des positions
de ses chauffeurs, tout en donnant à ces derniers un moyen direct de signaler
leur disponibilité ?

\section{Besoins auxquels répond la solution}

\begin{itemize}
  \item permettre au chauffeur de saisir son accès, d'autoriser la localisation
        et de passer en ligne ou hors ligne ;
  \item transmettre périodiquement une position géographique quand le suivi est
        actif ;
  \item donner à l'administrateur une carte, une liste de chauffeurs et des
        indicateurs de fraîcheur et de disponibilité ;
  \item actualiser l'affichage lors de la réception de nouveaux événements,
        sans dépendre uniquement d'une consultation manuelle ;
  \item permettre la création de fiches chauffeur depuis le tableau de bord ;
  \item séparer les données et les accès selon le centre, et distinguer
        l'administration d'un centre de l'administration de la plateforme.
\end{itemize}

\section{Objectifs et périmètre}

L'objectif de la version actuelle est de fournir un \important{socle
opérationnel de suivi GPS}. Le dépôt documente un échange réel entre une
application mobile et le tableau de bord. La position est suivie lorsque
l'application chauffeur est au premier plan ; l'envoi est réalisé à intervalle
rapproché, nominalement autour de sept secondes, selon la disponibilité du
service de localisation.

Le périmètre actuel couvre la gestion des chauffeurs et centres, la disponibilité,
l'enregistrement des positions et leur affichage en direct. La répartition des
courses, la gestion complète des véhicules, la facturation, les flux de travail
ERP, le suivi en arrière-plan et l'apprentissage automatique sont hors du
périmètre livré à ce stade.

\chapter{Solution fonctionnelle}

\section{Acteurs et responsabilités}

\begin{description}[style=nextline,leftmargin=3.5cm,labelwidth=3.2cm]
  \item[Chauffeur] utilise l'application mobile, renseigne son jeton d'accès,
        autorise la localisation et indique s'il est en service.
  \item[Administrateur de centre] se connecte au tableau de bord de son centre,
        suit l'état et la position des chauffeurs et peut créer une fiche
        chauffeur.
  \item[Responsable de plateforme] utilise une console séparée pour consulter
        les centres, en créer un et consulter le journal d'audit de la
        plateforme.
\end{description}

\section{Principaux parcours}

\subsection{Suivi d'un chauffeur}

Le chauffeur active son statut en ligne depuis l'application. Après
l'autorisation du système mobile, l'application obtient des coordonnées et les
transmet à l'API avec son identifiant d'accès. L'API vérifie l'accès, valide les
valeurs reçues et écarte les positions dont la précision dépasse le seuil
prévu. Les positions acceptées sont enregistrées puis diffusées au tableau de
bord du centre concerné.

Le tableau de bord met à jour la carte et les informations du chauffeur. Il
distingue notamment l'état en ligne, l'état hors ligne et une position devenue
ancienne. Une position récente est considérée comme en ligne par l'interface
lorsqu'elle date de moins d'une minute et que le statut du chauffeur est actif.

\subsection{Administration du centre}

Après authentification, l'administrateur consulte la carte et la liste des
chauffeurs de son centre. Il peut rechercher un chauffeur et créer une nouvelle
fiche. Un jeton d'accès associé au chauffeur est présenté une seule fois au
moment de sa création afin de pouvoir être transmis au chauffeur.

\subsection{Administration de la plateforme}

Le responsable de plateforme accède à une console distincte de celle des
centres. Dans l'état documenté, elle permet de consulter la liste des centres
et leurs principaux compteurs, d'ajouter un centre et de consulter le journal
d'audit associé à la console.

\section{Diagramme de cas d'utilisation}

\begin{figure}[htbp]
\centering
\resizebox{0.96\linewidth}{!}{%
\begin{tikzpicture}[
  actorline/.style={draw=fleetgray,line width=0.8pt},
  usecase/.style={draw=fleetblue,fill=blue!4,ellipse,minimum width=4.2cm,
                  minimum height=0.85cm,align=center},
  boundary/.style={draw=fleetgreen,line width=1pt,rounded corners=2mm}
]
  \draw[boundary] (2.1,-3.55) rectangle (11.4,4.15);
  \node[anchor=north west,text=fleetgreen,font=\bfseries] at (2.35,3.98)
    {Fleet Tracker};

  \node[align=center] (driver) at (0,1.0) {\textbf{Chauffeur}};
  \draw (0,2.45) circle (0.16);
  \draw (0,2.29) -- (0,1.75);
  \draw (0,2.12) -- (-0.22,1.91);
  \draw (0,2.12) -- (0.22,1.91);
  \draw (0,1.75) -- (-0.18,1.48);
  \draw (0,1.75) -- (0.18,1.48);

  \node[align=center] (admin) at (0,-0.2) {\textbf{Administrateur}\\\textbf{de centre}};
  \draw (0,-0.58) circle (0.16);
  \draw (0,-0.74) -- (0,-1.28);
  \draw (0,-0.91) -- (-0.22,-1.12);
  \draw (0,-0.91) -- (0.22,-1.12);
  \draw (0,-1.28) -- (-0.18,-1.55);
  \draw (0,-1.28) -- (0.18,-1.55);

  \node[align=center] (owner) at (13.4,-3.25) {\textbf{Responsable de}\\\textbf{plateforme}};
  \draw (13.4,-1.88) circle (0.16);
  \draw (13.4,-2.04) -- (13.4,-2.58);
  \draw (13.4,-2.21) -- (13.18,-2.42);
  \draw (13.4,-2.21) -- (13.62,-2.42);
  \draw (13.4,-2.58) -- (13.22,-2.85);
  \draw (13.4,-2.58) -- (13.58,-2.85);

  \node[usecase] (availability) at (5.05,2.75) {Indiquer sa disponibilité};
  \node[usecase] (sendloc) at (8.8,2.75) {Transmettre sa position};
  \node[usecase] (adminlogin) at (5.05,1.15) {S'authentifier au centre};
  \node[usecase] (monitor) at (8.8,1.15) {Suivre les chauffeurs\\sur la carte et la liste};
  \node[usecase] (create) at (5.05,-0.55) {Créer une fiche chauffeur};
  \node[usecase] (ownerlogin) at (8.8,-0.55) {S'authentifier à la plateforme};
  \node[usecase] (centers) at (5.05,-2.25) {Consulter et créer des centres};
  \node[usecase] (audit) at (8.8,-2.25) {Consulter le journal d'audit};

  \draw[actorline] (0.18,2.15) -- (availability.west);
  \draw[actorline] (0.18,2.05) -- (sendloc.west);
  \draw[actorline] (0.18,-0.55) -- (adminlogin.west);
  \draw[actorline] (0.18,-0.65) -- (monitor.west);
  \draw[actorline] (0.18,-0.75) -- (create.west);
  \draw[actorline] (13.18,-2.05) -- (ownerlogin.east);
  \draw[actorline] (13.18,-2.17) -- (centers.east);
  \draw[actorline] (13.18,-2.30) -- (audit.east);
\end{tikzpicture}%
}
\caption{Acteurs et fonctions principales du système.}
\label{fig:use-cases}
\end{figure}

\chapter{Architecture du système}

\section{Vue d'ensemble}

L'architecture sépare les interfaces, le traitement métier et la persistance.
L'application mobile et les deux interfaces web échangent avec un même backend.
Celui-ci centralise les règles de validation, l'authentification, le filtrage
des données par centre, l'enregistrement des positions et la diffusion des
événements. PostgreSQL conserve les données relationnelles et l'historique des
positions.

\begin{figure}[htbp]
\centering
\resizebox{\linewidth}{!}{%
\begin{tikzpicture}[
  box/.style={draw=fleetblue,fill=blue!4,rounded corners=2mm,
              text width=3.5cm,minimum height=1.15cm,align=center},
  external/.style={draw=fleetgray,fill=gray!7,rounded corners=2mm,
                   text width=2.8cm,minimum height=0.9cm,align=center},
  flow/.style={-{Stealth[length=2.5mm]},draw=fleetgreen,line width=1pt},
  returnflow/.style={-{Stealth[length=2.5mm]},draw=fleetblue,line width=1pt},
  label/.style={font=\small,fill=white,inner sep=2pt,align=center}
]
  \node[box] (mobile) at (-5,2.3)
    {\textbf{Application chauffeur}\\React Native / TypeScript};
  \node[box] (dashboard) at (-5,-0.2)
    {\textbf{Tableau de bord du centre}\\Navigateur, carte et liste};
  \node[box] (platform) at (-5,-2.7)
    {\textbf{Console de plateforme}\\Gestion des centres et audit};
  \node[box,text width=4.4cm,minimum height=1.55cm] (api) at (1,0)
    {\textbf{API centrale}\\FastAPI\\Validation, accès et diffusion temps réel};
  \node[box] (db) at (7,1.5)
    {\textbf{Base relationnelle}\\PostgreSQL};
  \node[external] (tiles) at (7,-1.5)
    {\textbf{Fond cartographique}\\Tuiles OpenStreetMap};
  \node[external] (gps) at (-9,2.3)
    {\textbf{Services de localisation}\\du téléphone};

  \draw[flow] (gps) -- node[label,above] {Position} (mobile);
  \draw[flow] (mobile.east) -- node[label,above] {Requêtes HTTP\\(statut et positions)} (api.west);
  \draw[returnflow] (api.west) -- node[label,below] {Confirmation} (mobile.east);
  \draw[flow] (dashboard.east) -- node[label,above] {HTTP : consultation\\et administration} (api.west);
  \draw[returnflow] (api.west) -- node[label,below] {WebSocket :\\événements en direct} (dashboard.east);
  \draw[flow] (platform.east) -- node[label,above] {HTTP : centres\\et audit} (api.west);
  \draw[flow] (api.east) -- node[label,above] {Lecture et écriture\\asynchrones} (db.west);
  \draw[flow] (dashboard.south east) -- node[label,below,sloped] {Chargement du fond\\de carte} (tiles.west);
\end{tikzpicture}%
}
\caption{Architecture logique et principaux échanges.}
\label{fig:architecture}
\end{figure}

\section{Composants}

\begin{description}[style=nextline,leftmargin=3.6cm,labelwidth=3.3cm]
  \item[Application chauffeur] application mobile Android et iOS. Elle gère
        l'accès du chauffeur, le statut de service et l'envoi des positions
        autorisées.
  \item[Tableau de bord web] interface de supervision destinée aux
        administrateurs de centre. Elle affiche une carte, des états
        synthétiques, une liste et des fonctions de gestion des chauffeurs.
  \item[Console de plateforme] interface web séparée pour les fonctions
        transversales portant sur les centres et le journal d'audit.
  \item[Backend] API centrale Python. Elle valide les requêtes, vérifie les
        accès, applique le périmètre du centre, enregistre les données et
        relaie les événements de suivi.
  \item[Base de données] stockage relationnel PostgreSQL pour les comptes, les
        centres, les chauffeurs, les positions et les journaux.
  \item[Services externes] service de localisation du téléphone pour les
        coordonnées, et tuiles OpenStreetMap comme fond de carte.
\end{description}

\section{Flux temps réel}

Les opérations de gestion et de collecte utilisent une API HTTP. Lorsqu'une
position acceptée est enregistrée, le backend peut diffuser un événement par
WebSocket aux tableaux de bord authentifiés du même centre. Le navigateur met
alors à jour la position sans attendre le prochain rafraîchissement complet.
La liste reste également rechargeable afin de récupérer un état initial ou de
se resynchroniser.

\chapter{Données et modèle conceptuel}

\section{Organisation des données}

La base de données est organisée autour du centre, qui constitue le périmètre
de gestion des chauffeurs, administrateurs et zones. Une position est liée à
un chauffeur et à son centre ; les données sont horodatées afin de distinguer
la dernière position de l'historique. Des tables séparées conservent les
sessions d'administration et les événements d'audit.

Le schéma ci-dessous est une représentation conceptuelle simplifiée du modèle
présent dans le projet. Les clés primaires et étrangères sont indiquées par
\textbf{PK} et \textbf{FK}. Les cardinalités décrivent les associations
principales, sans détailler les contraintes et tous les champs de stockage.

\begin{landscape}
\begin{figure}[p]
\centering
\resizebox{0.98\linewidth}{!}{%
\begin{tikzpicture}[
  entity/.style={draw=fleetblue,fill=blue!3,rounded corners=1mm,
                 text width=2.8cm,minimum height=0.9cm,align=center,
                 font=\small},
  relation/.style={draw=fleetgray,line width=0.8pt},
  relabel/.style={font=\scriptsize,fill=white,inner sep=1.5pt}
]
  \node[entity] (center) at (0,1.3)
    {\textbf{Centre}\\Identifiant (PK)\\Nom, ville, fuseau horaire};
  \node[entity] (admin) at (3.7,3.4)
    {\textbf{Administrateur}\\Identifiant (PK)\\Centre (FK), courriel, état};
  \node[entity] (driver) at (3.7,1.3)
    {\textbf{Chauffeur}\\Identifiant (PK)\\Centre (FK), identité, statut};
  \node[entity] (zone) at (3.7,-0.8)
    {\textbf{Zone}\\Identifiant (PK)\\Centre (FK), nom, limites};
  \node[entity] (adminsession) at (7.3,3.4)
    {\textbf{Session admin}\\Identifiant (PK)\\Administrateur (FK), échéance};
  \node[entity] (audit) at (7.3,1.3)
    {\textbf{Journal du centre}\\Identifiant (PK)\\Centre et admin (FK), action};
  \node[entity] (vehicle) at (7.3,-0.8)
    {\textbf{Véhicule}\\Identifiant (PK)\\Chauffeur (FK), immatriculation};
  \node[entity] (location) at (10.9,2.4)
    {\textbf{Position chauffeur}\\Identifiant (PK)\\Chauffeur et centre (FK),\\coordonnées, date};
  \node[entity] (statuslog) at (10.9,0.2)
    {\textbf{Historique de statut}\\Identifiant (PK)\\Chauffeur (FK), statut, date};
  \node[entity] (superadmin) at (14.5,3.4)
    {\textbf{Responsable plateforme}\\Identifiant (PK)\\Courriel, état};
  \node[entity] (platformsession) at (18.1,3.4)
    {\textbf{Session plateforme}\\Identifiant (PK)\\Responsable (FK), échéance};
  \node[entity] (platformaudit) at (18.1,1.0)
    {\textbf{Journal plateforme}\\Identifiant (PK)\\Responsable (FK), événement};

  \draw[relation] (center) -- node[relabel,above,sloped] {1 à plusieurs} (admin);
  \draw[relation] (center) -- node[relabel,above] {1 à plusieurs} (driver);
  \draw[relation] (center) -- node[relabel,below,sloped] {1 à plusieurs} (zone);
  \draw[relation] (admin) -- node[relabel,above] {1 à plusieurs} (adminsession);
  \draw[relation] (admin) -- node[relabel,above,sloped] {acteur} (audit);
  \draw[relation] (driver) -- node[relabel,above,sloped] {0 ou 1} (vehicle);
  \draw[relation] (driver) -- node[relabel,above,sloped] {1 à plusieurs} (location);
  \draw[relation] (driver) -- node[relabel,below,sloped] {1 à plusieurs} (statuslog);
  \draw[relation] (superadmin) -- node[relabel,above] {1 à plusieurs} (platformsession);
  \draw[relation] (superadmin) -- node[relabel,above,sloped] {1 à plusieurs} (platformaudit);
\end{tikzpicture}%
}
\caption{Diagramme conceptuel des principales entités persistées.}
\label{fig:database}
\end{figure}
\end{landscape}

\section{Lecture du modèle}

\begin{itemize}
  \item Un centre regroupe ses administrateurs, ses chauffeurs et ses zones.
  \item Un chauffeur peut être associé à un véhicule et possède plusieurs
        positions enregistrées ainsi qu'un historique de statut.
  \item Les administrateurs disposent de sessions et leurs opérations de
        gestion peuvent être consignées dans le journal du centre.
  \item Les responsables de plateforme possèdent des sessions distinctes et
        leurs opérations sont consignées dans un journal séparé.
\end{itemize}

\chapter{Séquence de suivi de la position}

\section{Scénario principal}

Le diagramme ci-dessous détaille un cycle représentatif : le chauffeur passe
en ligne, l'application obtient une position, l'API contrôle et enregistre
l'information, puis le tableau de bord reçoit l'événement. Les positions
signalées avec une précision supérieure à cinquante mètres sont ignorées par
le backend. Le suivi actuel est limité à l'utilisation de l'application au
premier plan.

\begin{figure}[p]
\centering
\resizebox{0.98\linewidth}{!}{%
\begin{tikzpicture}[
  participant/.style={draw=fleetblue,fill=blue!4,rounded corners=1mm,
                      text width=2.15cm,minimum height=0.9cm,align=center,
                      font=\small},
  message/.style={-{Stealth[length=2.4mm]},draw=fleetgreen,line width=0.9pt},
  response/.style={{Stealth[length=2.4mm]}-,draw=fleetblue,line width=0.9pt},
  event/.style={font=\scriptsize,fill=white,inner sep=2pt,align=center},
  lifeline/.style={draw=fleetgray,densely dashed}
]
  \foreach \id/\x/\name in {
    A/0/{Chauffeur},
    B/2.7/{Application\\mobile},
    C/5.4/{Service GPS\\du téléphone},
    D/8.1/{API\\centrale},
    E/10.8/{Base de\\données},
    F/13.5/{Tableau de\\bord}
  } {
    \node[participant] (p\id) at (\x,0) {\name};
    \draw[lifeline] (\x,-0.55) -- (\x,-9.55);
  }

  \draw[message] ([yshift=-0.9cm]pA.south) -- node[event,above] {1. Active le statut\\en ligne} ([yshift=-0.9cm]pB.south);
  \draw[message] ([yshift=-1.45cm]pB.south) -- node[event,above] {2. Demande l'accès\\à la localisation} ([yshift=-1.45cm]pC.south);
  \draw[response] ([yshift=-2cm]pC.south) -- node[event,above] {3. Autorisation\\accordée} ([yshift=-2cm]pB.south);
  \draw[message] ([yshift=-2.55cm]pB.south) -- node[event,above] {4. Envoie le statut et\\l'identifiant d'accès} ([yshift=-2.55cm]pD.south);
  \draw[message] ([yshift=-3.1cm]pD.south) -- node[event,above] {5. Vérifie l'accès et\\mémorise le statut} ([yshift=-3.1cm]pE.south);
  \draw[response] ([yshift=-3.65cm]pE.south) -- node[event,above] {6. Confirmation} ([yshift=-3.65cm]pD.south);
  \draw[response] ([yshift=-4.2cm]pD.south) -- node[event,above] {7. Confirme l'état\\en ligne} ([yshift=-4.2cm]pB.south);
  \draw[message] ([yshift=-4.75cm]pB.south) -- node[event,above] {8. Demande une\\position actuelle} ([yshift=-4.75cm]pC.south);
  \draw[response] ([yshift=-5.3cm]pC.south) -- node[event,above] {9. Renvoie position,\\précision et date} ([yshift=-5.3cm]pB.south);
  \draw[message] ([yshift=-5.85cm]pB.south) -- node[event,above] {10. Transmet la position\\avec l'identifiant d'accès} ([yshift=-5.85cm]pD.south);
  \draw[message] ([yshift=-6.4cm]pD.south) -- node[event,above] {11. Valide puis enregistre\\la position acceptée} ([yshift=-6.4cm]pE.south);
  \draw[response] ([yshift=-6.95cm]pE.south) -- node[event,above] {12. Écriture confirmée} ([yshift=-6.95cm]pD.south);
  \draw[message] ([yshift=-7.5cm]pD.south) -- node[event,above] {13. Diffuse l'événement\\au centre concerné} ([yshift=-7.5cm]pF.south);
  \draw[message] ([yshift=-8.05cm]pF.south) -- ++(0.8,0) -- ++(0,-0.35)
    -- node[event,left] {14. Actualise la carte\\et l'état affiché} ++(-0.8,0);

  \node[draw=fleetgray,fill=gray!6,rounded corners=1mm,
        text width=8.2cm,align=center,font=\scriptsize]
        at (6.75,-9.1)
        {Si la précision reçue dépasse le seuil de 50 m, la position est écartée
        et ne produit pas de nouvel événement cartographique.};
\end{tikzpicture}%
}
\caption{Séquence simplifiée de transmission et d'affichage d'une position.}
\label{fig:sequence}
\end{figure}

\chapter{Choix technologiques}

\section{Principes de choix}

Les choix privilégient une séparation claire entre les interfaces et les
services, des technologies adaptées au suivi géographique et une base
relationnelle appropriée aux associations entre centres, comptes, chauffeurs
et historiques. Ils permettent également de construire progressivement le
produit sans faire dépendre les applications clientes l'une de l'autre.

\begin{table}[htbp]
\centering
\small
\begin{tabularx}{\linewidth}{>{\bfseries}p{3.2cm} X X}
\toprule
Composant & Technologie retenue & Justification synthétique \\
\midrule
API serveur & Python et FastAPI & Permet de structurer les échanges HTTP et
temps réel, de valider les données reçues et de centraliser les règles
fonctionnelles. \\
Persistance & PostgreSQL & Convient aux données structurées, aux relations
entre entités et à la conservation des historiques utiles au suivi. \\
Accès aux données & SQLAlchemy asynchrone et Alembic & Fournit une couche
relationnelle cohérente avec le traitement asynchrone et une évolution
contrôlée du schéma. \\
Application mobile & React Native et TypeScript & Permet de partager une
approche de développement entre Android et iOS tout en accédant aux capacités
de géolocalisation du téléphone. \\
Tableau de bord & HTML, CSS et JavaScript & Offre une interface navigateur
légère et indépendante du cycle de développement mobile. \\
Cartographie & Leaflet et OpenStreetMap & Associe une bibliothèque de carte
interactive à un fond cartographique adapté à la visualisation des positions. \\
Communication en direct & WebSocket & Permet au backend de pousser les nouveaux
événements de statut et de position au tableau de bord connecté. \\
\bottomrule
\end{tabularx}
\caption{Technologies majeures et raisons de leur choix.}
\label{tab:stack}
\end{table}

\section{Cohérence de l'ensemble}

L'API est le point commun des trois interfaces, ce qui évite de disperser la
validation et les règles d'accès. La base relationnelle conserve un historique
exploitable plutôt que de ne garder qu'une position en mémoire. Les échanges
HTTP couvrent les opérations classiques, tandis que le WebSocket répond au
besoin de mise à jour immédiate de la carte. Enfin, l'application mobile est
adaptée à l'usage du chauffeur sur le terrain.

\chapter{État actuel, limites et perspectives}

\section{État fonctionnel}

Le projet comprend les éléments suivants :
\begin{itemize}
  \item une API de suivi et de gestion des accès ;
  \item une application mobile chauffeur avec disponibilité et envoi de
        positions au premier plan ;
  \item un tableau de bord web présentant chauffeurs, carte et états de
        fraîcheur ;
  \item une console séparée pour la gestion des centres et la consultation
        d'un journal d'audit de plateforme ;
  \item un stockage PostgreSQL structuré par centre et des migrations de
        schéma.
\end{itemize}

La documentation du projet fait état d'un parcours de bout en bout vérifié :
une position réelle provenant d'un téléphone ou d'un émulateur apparaît sur
le tableau de bord.

\section{Limites identifiées}

\begin{itemize}
  \item Le suivi de localisation s'arrête lorsque l'application quitte son
        écran principal ou n'est plus au premier plan.
  \item Le système n'attribue pas de course et ne gère ni la réservation, ni
        l'affectation, ni la facturation.
  \item Les fonctionnalités ERP et d'aide à la décision par apprentissage
        automatique restent des perspectives et non des fonctions livrées.
  \item Le dépôt ne décrit pas encore de déploiement cloud ni de chaîne
        d'intégration et de livraison continue.
  \item La couverture des tests est encore en cours d'élargissement par rapport
        aux parcours fonctionnels complets.
\end{itemize}

\section{Perspectives possibles}

Les prolongements naturels peuvent inclure un suivi en arrière-plan soumis aux
autorisations mobiles, une gestion plus complète des véhicules, des parcours
de dispatching et des écrans de rapports. Ils devraient être introduits
progressivement, en s'appuyant sur les besoins des utilisateurs, les règles de
protection des données de localisation et les capacités de l'infrastructure
cible.

\chapter{Conclusion}

Fleet Tracker apporte une réponse concrète au besoin de visibilité sur une
flotte de taxis : le chauffeur signale sa disponibilité et transmet sa
position, tandis que le responsable suit l'activité depuis une interface
cartographique. Le backend et la base de données assurent le lien central entre
ces usages, avec une séparation des données par centre et une mise à jour en
direct du tableau de bord.

La version actuelle constitue une fondation fonctionnelle de suivi GPS, et
non un système complet de dispatching ou de gestion de transport. Son
architecture modulaire permet d'envisager ces évolutions sans confondre les
capacités déjà présentes et les objectifs futurs.

\appendix
\chapter{Glossaire}

\begin{description}[style=nextline,leftmargin=3.2cm,labelwidth=2.9cm]
  \item[API] Interface par laquelle les applications échangent des données
        avec le service backend.
  \item[Centre] Périmètre organisationnel auquel sont rattachés des
        administrateurs, des chauffeurs et des données de suivi.
  \item[Position] Coordonnées géographiques horodatées, éventuellement
        complétées par une vitesse, une direction et une précision.
  \item[WebSocket] Mécanisme de communication qui permet au serveur de pousser
        des événements vers une interface déjà connectée.
  \item[Fraîcheur] Âge de la dernière position reçue, utilisé pour distinguer
        une information récente d'une information devenue ancienne.
\end{description}

\end{document}
